\documentclass[12pt]{article}
\usepackage{amsmath,amssymb,amsthm}

\newtheorem{theorem}{Theorem}
\newtheorem{lemma}[theorem]{Lemma}
\newtheorem{corollary}[theorem]{Corollary}

\begin{document}

\section*{Generic 2-Vertex Cut Pullback Lemma}

This document establishes a unified pullback mechanism for 2-vertex cuts in graphs generated by reduction gadgets in the class $\mathcal{Q}$.

\begin{lemma}[Connectedness of $G \setminus H$]
\label{lem:isolated_neighborhood}
Let $G \in \mathcal{Q}$ contain a certified occurrence $H$ with boundary $B$ where $|B|=4$ and the only edges between $V(H)$ and $V(G)\setminus V(H)$ are the four edges from $B$ into $H$.
\end{lemma}
\begin{proof}
By 3-connectivity of $G$, any component $C$ of $G[V(G)\setminus V(H)]$ satisfies $|E(C, V(H))| \ge 3$. If at least two such components existed, the total edges leaving $V(G)\setminus V(H)$ would be $\ge 6$, contradicting the certified-occurrence definition of exactly 4 boundary edges.
\end{proof}

\begin{lemma}[Generic 2-Vertex Cut Pullback]
\label{lem:pullback_generic}
Let $G \in \mathcal{Q}$ contain a certified occurrence $H$ with boundary $B$. Let $G'$ be the reduced graph where $H$ is replaced by a gadget $H'$ consisting of exactly two adjacent vertices $x, y$. Let $B_x, B_y$ be a partition of $B$ into two disjoint sets of size 2 such that $N_{G'}(x) \setminus \{y\} = B_x$ and $N_{G'}(y) \setminus \{x\} = B_y$.

If $G'$ has a 2-vertex cut $A = \{p, q\}$, then either:
\begin{enumerate}
    \item[(i)] $A \subseteq V(G) \setminus V(H)$ and $A$ is a 2-vertex cut of $G$; or
    \item[(ii)] exactly one of $\{p, q\}$ is a gadget vertex ($p \in \{x, y\}$) and $q \in V(G) \setminus V(H)$, and there exists a nonempty vertex set $K \subseteq V(G) \setminus V(H)$ such that every edge leaving $K$ in $G$ has its endpoint in $\{q\} \cup B_p$ (where $B_p$ is the set of terminals adjacent to $p$).
\end{enumerate}
\end{lemma}

\begin{proof}
Let $A = \{p, q\}$ be a 2-vertex cut in $G'$.

\medskip
\noindent\textbf{Case 1: $p, q \notin \{x, y\}$.} \\
Since $xy \in E(G')$, the vertices $x$ and $y$ lie in the same component of $G' - A$; denote it $C_0$. Since $A$ is a cut, there exists another component $C_1$.
Then $V(C_1) \subseteq V(G) \setminus V(H)$. Since any terminal $b \in B \setminus A$ is adjacent to $x \in C_0$ or $y \in C_0$, we have $B \subseteq V(C_0) \cup A$.
Restoring $H$ adds edges only between $V(H)$ and $B$. Thus $C_1$ remains disjoint from $V(H)$ and has no edges to $V(H)$ in $G$. Since edges in $G \setminus H$ are identical to those in $G' \setminus H'$, all edges leaving $C_1$ in $G$ must land in $A$. Thus $A$ is a 2-vertex cut in $G$.

\medskip
\noindent\textbf{Case 2: $\{p, q\} = \{x, y\}$.} \\
$G' - \{x, y\} = G[\,V(G) \setminus V(H)\,]$, which is connected by Lemma~\ref{lem:isolated_neighborhood}. Thus $\{x, y\}$ cannot be a 2-vertex cut.

\medskip
\noindent\textbf{Case 3: Exactly one of $\{p, q\}$ lies in $\{x, y\}$.} \\
Assume $p = x$ and $q \in V(G) \setminus V(H)$. Let $C_0$ be the component of $G' - \{x, q\}$ containing $y$. Let $C_1$ be another component.
Define $K = V(C_1) \setminus B_x$. Since $C_1$ is disjoint from $C_0$ (which contains $y$), $C_1 \subseteq V(G) \setminus V(H)$.
Because $y \in C_0$, its non-gadget neighbors $B_y \setminus \{q\}$ lie in $C_0 \cup \{q\}$ in $G' - \{x, q\}$. Thus $C_1 \cap B_y = \emptyset$, implying $C_1 \cap B \subseteq B_x$.
It follows that $K = V(C_1) \setminus B$.

If $V(C_1) \nsubseteq B_x$, then $K \neq \emptyset$ immediately. Otherwise pick $b \in V(C_1) \cap B_x$.
In $G'$, the vertex $b \neq q$ has degree $3$ with neighbors $\{x\} \cup N$, where $N \subseteq V(G)\setminus V(H)$ consists of two distinct non-gadget neighbors.
After deleting $\{x,q\}$, at most one vertex of $N$ can equal $q$, so there exists $w \in N \setminus \{q\}$.
Then $bw$ is an edge of $G' - \{x,q\}$, hence $w \in V(C_1)$. Moreover $w \notin B_x$ (since $b$ has at most one neighbor in the 2-set $B_x$, namely possibly the other element of $B_x$).
Therefore $w \in V(C_1)\setminus B_x = K$, so $K \neq \emptyset$.

Since $K \cap B = \emptyset$, $K$ has no neighbors in $V(H)$ in $G$. In $G'$, any edge from $v \in K$ to $w \notin K \cup \{q\} \cup B_x$ would imply $w$ is a neighbor of $C_1$ in $G' - \{x, q\}$ outside $C_1 \cup \{q\}$, which is impossible as $x$ is the only other potential neighbor and $x$ only connects to $B_x$. Thus neighbors of $K$ in $G$ are contained in $\{q\} \cup B_x$.
\end{proof}

\begin{corollary}[Configuration Instantiations]
The generic lemma holds for each configuration in the manuscript under the following terminal partitions:
\begin{itemize}
    \item \textbf{$C_2$}: $B_x = \{u_1, u_6\}$, $B_y = \{u_4, u_5\}$.
    \item \textbf{Refined $C_4$}: $B_x = \{u_1, u_3\}$, $B_y = \{u_2, u_4\}$.
    \item \textbf{$C_P$}: $B_x = \{r, s\}$, $B_y = \{u_2, u_4\}$.
\end{itemize}
\end{corollary}

\end{document}
